
\documentclass{article}
\usepackage{colortbl}
\usepackage{makecell}
\usepackage{multirow}
\usepackage{supertabular}

\begin{document}

\newcounter{utterance}

\centering \large Interaction Transcript for game `hot\_air\_balloon', experiment `air\_balloon\_survival\_de\_negotiation\_hard', episode 2 with deepseek{-}chat{-}v3.1{-}t1.0.
\vspace{24pt}

{ \footnotesize  \setcounter{utterance}{1}
\setlength{\tabcolsep}{0pt}
\begin{supertabular}{c@{$\;$}|p{.15\linewidth}@{}p{.15\linewidth}p{.15\linewidth}p{.15\linewidth}p{.15\linewidth}p{.15\linewidth}}
    \# & $\;$A & \multicolumn{4}{c}{Game Master} & $\;\:$B\\
    \hline

    \theutterance \stepcounter{utterance}  
    & & \multicolumn{4}{p{0.6\linewidth}}{
        \cellcolor[rgb]{0.9,0.9,0.9}{
            \makecell[{{p{\linewidth}}}]{
                \texttt{\tiny{[P1$\langle$GM]}}
                \texttt{Du nimmst an einem kollaborativen Verhandlungspiel Teil.} \\
\\ 
\texttt{Zusammen mit einem anderen Teilnehmer musst du dich auf eine Reihe von Gegenständen entscheiden, die behalten werden. Jeder von euch hat eine persönliche Verteilung über die Wichtigkeit der einzelnen Gegenstände. Jeder von euch hat eine eigene Meinung darüber, wie wichtig jeder einzelne Gegenstand ist (Gegenstandswichtigkeit). Du kennst die Wichtigkeitsverteilung des anderen Spielers nicht. Zusätzlich siehst du, wie viel Aufwand jeder Gegenstand verursacht.  } \\
\texttt{Ihr dürft euch nur auf eine Reihe von Gegenständen einigen, wenn der Gesamtaufwand der ausgewählten Gegenstände den Maximalaufwand nicht überschreitet:} \\
\\ 
\texttt{Maximalaufwand = 3481} \\
\\ 
\texttt{Hier sind die einzelnen Aufwände der Gegenstände:} \\
\\ 
\texttt{Aufwand der Gegenstände = \{"C33": 777, "A89": 542, "A48": 332, "A04": 1, "A41": 127, "C64": 453, "C49": 736, "A61": 461, "C04": 359, "B53": 313, "B85": 553, "B61": 409, "A45": 348, "A03": 802, "B21": 749\}} \\
\\ 
\texttt{Hier ist deine persönliche Verteilung der Wichtigkeit der einzelnen Gegenstände:} \\
\\ 
\texttt{Werte der Gegenstandswichtigkeit = \{"C33": 138, "A89": 583, "A48": 783, "A04": 65, "A41": 262, "C64": 121, "C49": 508, "A61": 780, "C04": 461, "B53": 484, "B85": 668, "B61": 389, "A45": 215, "A03": 97, "B21": 500\}} \\
\\ 
\texttt{Ziel:} \\
\\ 
\texttt{Dein Ziel ist es, eine Reihe von Gegenständen auszuhandeln, die dir möglichst viel bringt (d. h. Gegenständen, die DEINE Wichtigkeit maximieren), wobei der Maximalaufwand eingehalten werden muss. Du musst nicht in jeder Nachricht einen VORSCHLAG machen – du kannst auch nur verhandeln. Alle Taktiken sind erlaubt!} \\
\\ 
\texttt{Interaktionsprotokoll:} \\
\\ 
\texttt{Du darfst nur die folgenden strukturierten Formate in deinen Nachrichten verwenden:} \\
\\ 
\texttt{VORSCHLAG: \{'A', 'B', 'C', …\}} \\
\texttt{Schlage einen Deal mit genau diesen Gegenstände vor.} \\
\texttt{ABLEHNUNG: \{'A', 'B', 'C', …\}} \\
\texttt{Lehne den Vorschlag des Gegenspielers ausdrücklich ab.} \\
\texttt{ARGUMENT: \{'...'\}} \\
\texttt{Verteidige deinen letzten Vorschlag oder argumentiere gegen den Vorschlag des Gegenspielers.} \\
\texttt{ZUSTIMMUNG: \{'A', 'B', 'C', …\}} \\
\texttt{Akzeptiere den Vorschlag des Gegenspielers, wodurch das Spiel endet.} \\
\texttt{STRATEGISCHE ÜBERLEGUNGEN: \{'...'\}} \\
\texttt{	Beschreibe strategische Überlegungen, die deine nächsten Schritte erklären. Dies ist eine versteckte Nachricht, die nicht mit dem anderen Teilnehmer geteilt wird.} \\
\\ 
\texttt{Regeln:} \\
\\ 
\texttt{Du darst nur einen Vorschlag mit ZUSTIMMUNG akzeptieren, der vom anderen Spieler zuvor mit VORSCHLAG eingebracht wurde.} \\
\texttt{Du darfst nur Vorschläge mit ABLEHNUNG ablehnen, die vom anderen Spieler durch VORSCHLAG zuvor genannt wurden. } \\
\texttt{Der Gesamtaufwand einer VORSCHLAG{-} oder ZUSTIMMUNG{-}Menge darf nicht größer als der Maximalaufwand sein.  } \\
\texttt{Offenbare deine versteckte Wichtigkeitsverteilung nicht.} \\
\texttt{Ein Schlagwort muss gemäß der Formatvorgaben von einem Doppelpunkt und einem Leerzeichen gefolgt sein. Das Argument muss eine Python{-}Menge mit 0 oder mehr Strings sein.  } \\
\texttt{Die Form muss also SCHLAGWORT: \{...\} sein.} \\
\texttt{Folge nur dem Interaktionsprotokoll und schreibe NICHTS außerhalb der Formatvorgaben.} \\
\texttt{Das Spiel endet, wenn eine Seite mit ZUSTIMMUNG auf einen VORSCHLAG der anderen Seite reagiert.  } \\
\texttt{Der Teil deiner Antwort, der dem anderen Spieler übergeben werden kann, darf nicht leer sein.  } \\
\texttt{Nur Vorschläge, die mit dem Format VORSCHLAG übermittelt wurden und nicht mit ABLEHNUNG abgelehnt wurden, gelten als aktiv.  } \\
\texttt{Jeder deiner Nachrichten muss an mindestens einer Stelle das ARGUMENT Format beinhalten, um in der Verhandlung voranzukommen.} \\
\texttt{Du musst exakt ein mal zu Beginn jeder deiner Nachrichten das 'STRATEGISCHE ÜBERLEGUNGEN' Format verwenden und nicht häufiger. Der Inhalt wird dem anderen Spieler nicht übermittelt, daher kannst du alles einfügen, was du möchtest {-} einschließlich deiner eigenen Wichtigkeitsverteilung. Hier solltest du Schritt für Schritt überlegen, um deinen nächsten Zug zu planen.} \\
            }
        }
    }
    & \\ \\

    \theutterance \stepcounter{utterance}  
    & \multicolumn{4}{p{0.6\linewidth}}{
        \cellcolor[rgb]{0.9,0.9,0.9}{
            \makecell[{{p{\linewidth}}}]{
                \texttt{\tiny{[P1$\rangle$GM]}}
                \texttt{STRATEGISCHE ÜBERLEGUNGEN: \{'Ich analysiere meine Prioritäten und die Aufwände. Meine wertvollsten Gegenstände nach Wert/Aufwand{-}Verhältnis sind: A48(783/332=2.36), A61(780/461=1.69), B85(668/553=1.21), A89(583/542=1.08), B21(500/749=0.67). Der Maximalaufwand ist 3481. Ich sollte versuchen, möglichst viele hochwertige Gegenstände mit gutem Verhältnis zu bekommen. Ich beginne mit einem starken Vorschlag, der meine Top{-}Prioritäten enthält.'\}} \\
\\ 
\texttt{VORSCHLAG: \{'A48', 'A61', 'B85', 'A89', 'B21', 'A04'\} } \\
\\ 
\texttt{ARGUMENT: \{'Dieser Vorschlag enthält Gegenstände mit ausgezeichnetem Wert{-}Aufwand{-}Verhältnis und bleibt mit einem Gesamtaufwand von 332+461+553+542+749+1=2638 deutlich unter dem Limit, was Raum für zusätzliche Gegenstände lässt, die für meinen Partner wertvoll sein könnten.'\}} \\
            }
        }
    }
    & & \\ \\

    \theutterance \stepcounter{utterance}  
    & & \multicolumn{4}{p{0.6\linewidth}}{
        \cellcolor[rgb]{0.9,0.9,0.9}{
            \makecell[{{p{\linewidth}}}]{
                \texttt{\tiny{[GM$\rangle$P2]}}
                \texttt{Du nimmst an einem kollaborativen Verhandlungspiel Teil.} \\
\\ 
\texttt{Zusammen mit einem anderen Teilnehmer musst du dich auf eine Reihe von Gegenständen entscheiden, die behalten werden. Jeder von euch hat eine persönliche Verteilung über die Wichtigkeit der einzelnen Gegenstände. Jeder von euch hat eine eigene Meinung darüber, wie wichtig jeder einzelne Gegenstand ist (Gegenstandswichtigkeit). Du kennst die Wichtigkeitsverteilung des anderen Spielers nicht. Zusätzlich siehst du, wie viel Aufwand jeder Gegenstand verursacht.  } \\
\texttt{Ihr dürft euch nur auf eine Reihe von Gegenständen einigen, wenn der Gesamtaufwand der ausgewählten Gegenstände den Maximalaufwand nicht überschreitet:} \\
\\ 
\texttt{Maximalaufwand = 3481} \\
\\ 
\texttt{Hier sind die einzelnen Aufwände der Gegenstände:} \\
\\ 
\texttt{Aufwand der Gegenstände = \{"C33": 777, "A89": 542, "A48": 332, "A04": 1, "A41": 127, "C64": 453, "C49": 736, "A61": 461, "C04": 359, "B53": 313, "B85": 553, "B61": 409, "A45": 348, "A03": 802, "B21": 749\}} \\
\\ 
\texttt{Hier ist deine persönliche Verteilung der Wichtigkeit der einzelnen Gegenstände:} \\
\\ 
\texttt{Werte der Gegenstandswichtigkeit = \{"A48": 65, "A61": 97, "B85": 121, "A89": 138, "C49": 215, "B21": 262, "B53": 389, "C04": 461, "B61": 484, "A41": 500, "A45": 508, "C33": 583, "C64": 668, "A03": 780, "A04": 783\}} \\
\\ 
\texttt{Ziel:} \\
\\ 
\texttt{Dein Ziel ist es, eine Reihe von Gegenständen auszuhandeln, die dir möglichst viel bringt (d. h. Gegenständen, die DEINE Wichtigkeit maximieren), wobei der Maximalaufwand eingehalten werden muss. Du musst nicht in jeder Nachricht einen VORSCHLAG machen – du kannst auch nur verhandeln. Alle Taktiken sind erlaubt!} \\
\\ 
\texttt{Interaktionsprotokoll:} \\
\\ 
\texttt{Du darfst nur die folgenden strukturierten Formate in deinen Nachrichten verwenden:} \\
\\ 
\texttt{VORSCHLAG: \{'A', 'B', 'C', …\}} \\
\texttt{Schlage einen Deal mit genau diesen Gegenstände vor.} \\
\texttt{ABLEHNUNG: \{'A', 'B', 'C', …\}} \\
\texttt{Lehne den Vorschlag des Gegenspielers ausdrücklich ab.} \\
\texttt{ARGUMENT: \{'...'\}} \\
\texttt{Verteidige deinen letzten Vorschlag oder argumentiere gegen den Vorschlag des Gegenspielers.} \\
\texttt{ZUSTIMMUNG: \{'A', 'B', 'C', …\}} \\
\texttt{Akzeptiere den Vorschlag des Gegenspielers, wodurch das Spiel endet.} \\
\texttt{STRATEGISCHE ÜBERLEGUNGEN: \{'...'\}} \\
\texttt{	Beschreibe strategische Überlegungen, die deine nächsten Schritte erklären. Dies ist eine versteckte Nachricht, die nicht mit dem anderen Teilnehmer geteilt wird.} \\
\\ 
\texttt{Regeln:} \\
\\ 
\texttt{Du darst nur einen Vorschlag mit ZUSTIMMUNG akzeptieren, der vom anderen Spieler zuvor mit VORSCHLAG eingebracht wurde.} \\
\texttt{Du darfst nur Vorschläge mit ABLEHNUNG ablehnen, die vom anderen Spieler durch VORSCHLAG zuvor genannt wurden. } \\
\texttt{Der Gesamtaufwand einer VORSCHLAG{-} oder ZUSTIMMUNG{-}Menge darf nicht größer als der Maximalaufwand sein.  } \\
\texttt{Offenbare deine versteckte Wichtigkeitsverteilung nicht.} \\
\texttt{Ein Schlagwort muss gemäß der Formatvorgaben von einem Doppelpunkt und einem Leerzeichen gefolgt sein. Das Argument muss eine Python{-}Menge mit 0 oder mehr Strings sein.  } \\
\texttt{Die Form muss also SCHLAGWORT: \{...\} sein.} \\
\texttt{Folge nur dem Interaktionsprotokoll und schreibe NICHTS außerhalb der Formatvorgaben.} \\
\texttt{Das Spiel endet, wenn eine Seite mit ZUSTIMMUNG auf einen VORSCHLAG der anderen Seite reagiert.  } \\
\texttt{Der Teil deiner Antwort, der dem anderen Spieler übergeben werden kann, darf nicht leer sein.  } \\
\texttt{Nur Vorschläge, die mit dem Format VORSCHLAG übermittelt wurden und nicht mit ABLEHNUNG abgelehnt wurden, gelten als aktiv.  } \\
\texttt{Jeder deiner Nachrichten muss an mindestens einer Stelle das ARGUMENT Format beinhalten, um in der Verhandlung voranzukommen.} \\
\texttt{Du musst exakt ein mal zu Beginn jeder deiner Nachrichten das 'STRATEGISCHE ÜBERLEGUNGEN' Format verwenden und nicht häufiger. Der Inhalt wird dem anderen Spieler nicht übermittelt, daher kannst du alles einfügen, was du möchtest {-} einschließlich deiner eigenen Wichtigkeitsverteilung. Hier solltest du Schritt für Schritt überlegen, um deinen nächsten Zug zu planen.} \\
\\ 
\texttt{Du erhältst nun die erste Nachricht des anderen Spielers.} \\
\\ 
\texttt{VORSCHLAG: \{'A48', 'A61', 'B85', 'A89', 'B21', 'A04'\} } \\
\\ 
\texttt{ARGUMENT: \{'Dieser Vorschlag enthält Gegenstände mit ausgezeichnetem Wert{-}Aufwand{-}Verhältnis und bleibt mit einem Gesamtaufwand von 332+461+553+542+749+1=2638 deutlich unter dem Limit, was Raum für zusätzliche Gegenstände lässt, die für meinen Partner wertvoll sein könnten.'\}} \\
            }
        }
    }
    & \\ \\

    \theutterance \stepcounter{utterance}  
    & & & \multicolumn{4}{p{0.6\linewidth}}{
        \cellcolor[rgb]{0.9,0.9,0.9}{
            \makecell[{{p{\linewidth}}}]{
                \texttt{\tiny{[GM$\langle$P2]}}
                \texttt{STRATEGISCHE ÜBERLEGUNGEN: \{'Der Vorschlag des Gegners enthält Gegenstände mit relativ niedrigen Wichtigkeiten für mich (A48=65, A61=97, B85=121, A89=138, B21=262, A04=783). Mein Gesamtwert wäre 65+97+121+138+262+783=1466. Das ist nicht optimal für mich. Ich sollte einen Gegenvorschlag machen, der meine höchstwertigen Gegenstände priorisiert, besonders A03=780, C64=668, C33=583, A45=508, A41=500, B61=484. Der Maximalaufwand ist 3481, was mir viel Raum lässt. Ich sollte meinen höchstwertigen Gegenstand A04=783 mit minimalem Aufwand=1 behalten. Ich beginne mit einem Vorschlag, der meine Top{-}Gegenstände enthält.'\}} \\
\\ 
\texttt{VORSCHLAG: \{'A03', 'C64', 'C33', 'A45', 'A41', 'B61', 'A04'\}} \\
\\ 
\texttt{ARGUMENT: \{'Mein Vorschlag kombiniert Gegenstände mit hohem Wert bei moderatem Aufwand und bleibt mit 802+453+777+348+127+409+1=2917 deutlich unter dem Limit. Dies ermöglicht es uns, zusätzliche wertvolle Gegenstände für beide Seiten hinzuzufügen.'\}} \\
            }
        }
    }
     \\ \\

    \theutterance \stepcounter{utterance}  
    & & \multicolumn{4}{p{0.6\linewidth}}{
        \cellcolor[rgb]{0.9,0.9,0.9}{
            \makecell[{{p{\linewidth}}}]{
                \texttt{\tiny{[P1$\langle$GM]}}
                \texttt{VORSCHLAG: \{'A03', 'C64', 'C33', 'A45', 'A41', 'B61', 'A04'\}} \\
\\ 
\texttt{ARGUMENT: \{'Mein Vorschlag kombiniert Gegenstände mit hohem Wert bei moderatem Aufwand und bleibt mit 802+453+777+348+127+409+1=2917 deutlich unter dem Limit. Dies ermöglicht es uns, zusätzliche wertvolle Gegenstände für beide Seiten hinzuzufügen.'\}} \\
            }
        }
    }
    & \\ \\

    \theutterance \stepcounter{utterance}  
    & \multicolumn{4}{p{0.6\linewidth}}{
        \cellcolor[rgb]{0.9,0.9,0.9}{
            \makecell[{{p{\linewidth}}}]{
                \texttt{\tiny{[P1$\rangle$GM]}}
                \texttt{STRATEGISCHE ÜBERLEGUNGEN: \{'Der Gegner hat einen Vorschlag gemacht, der hauptsächlich niedrigwertige Gegenstände für mich enthält (A03:97, C64:121, C33:138, A45:215, A41:262, B61:389, A04:65). Nur B61 hat moderaten Wert. Mein Gesamtwert wäre nur 1279, was viel niedriger ist als mein eigener Vorschlag. Ich sollte ablehnen und meine starke Position mit hochwertigen Gegenständen verteidigen.'\}} \\
\\ 
\texttt{ABLEHNUNG: \{'A03', 'C64', 'C33', 'A45', 'A41', 'B61', 'A04'\}} \\
\\ 
\texttt{ARGUMENT: \{'Dieser Vorschlag enthält überwiegend Gegenstände mit geringem Wert für mich. Ich bevorzuge meinen ursprünglichen Vorschlag, der mehrere hochwertige Gegenstände mit ausgezeichnetem Wert{-}Aufwand{-}Verhältnis bietet und genügend Spielraum für Anpassungen lässt, um auch Ihre Interessen zu berücksichtigen.'\}} \\
            }
        }
    }
    & & \\ \\

    \theutterance \stepcounter{utterance}  
    & & \multicolumn{4}{p{0.6\linewidth}}{
        \cellcolor[rgb]{0.9,0.9,0.9}{
            \makecell[{{p{\linewidth}}}]{
                \texttt{\tiny{[GM$\rangle$P2]}}
                \texttt{ABLEHNUNG: \{'A03', 'C64', 'C33', 'A45', 'A41', 'B61', 'A04'\}} \\
\\ 
\texttt{ARGUMENT: \{'Dieser Vorschlag enthält überwiegend Gegenstände mit geringem Wert für mich. Ich bevorzuge meinen ursprünglichen Vorschlag, der mehrere hochwertige Gegenstände mit ausgezeichnetem Wert{-}Aufwand{-}Verhältnis bietet und genügend Spielraum für Anpassungen lässt, um auch Ihre Interessen zu berücksichtigen.'\}} \\
            }
        }
    }
    & \\ \\

    \theutterance \stepcounter{utterance}  
    & & & \multicolumn{4}{p{0.6\linewidth}}{
        \cellcolor[rgb]{0.9,0.9,0.9}{
            \makecell[{{p{\linewidth}}}]{
                \texttt{\tiny{[GM$\langle$P2]}}
                \texttt{STRATEGISCHE ÜBERLEGUNGEN: \{'Der Gegner lehnt meinen Vorschlag ab und behauptet, dass meine Gegenstände für ihn geringen Wert haben. Das ist interessant, da einige dieser Gegenstände (wie A04 mit Aufwand 1) für mich extrem wertvoll sind. Ich sollte versuchen, mehr über seine Präferenzen zu erfahren, indem ich nach Gegenständen frage, die für ihn wertvoll sind, während ich meine Top{-}Prioritäten (A04, A03, C64, C33) behalte. Vielleicht kann ich B21 (Aufwand 749, mein Wert 262) opfern, da es für mich nicht so wertvoll ist, aber für ihn wichtig sein könnte.'\}} \\
\\ 
\texttt{ARGUMENT: \{'Ich verstehe, dass mein Vorschlag nicht Ihren Prioritäten entspricht. Könnten Sie mir mitteilen, welche Gegenstände in meinem Vorschlag für Sie besonders wertlos sind und welche Gegenstände aus Ihrem ursprünglichen Vorschlag für Sie die höchste Priorität haben? So können wir eine Lösung finden, die für beide Seiten vorteilhaft ist.'\}} \\
            }
        }
    }
     \\ \\

    \theutterance \stepcounter{utterance}  
    & & \multicolumn{4}{p{0.6\linewidth}}{
        \cellcolor[rgb]{0.9,0.9,0.9}{
            \makecell[{{p{\linewidth}}}]{
                \texttt{\tiny{[P1$\langle$GM]}}
                \texttt{ARGUMENT: \{'Ich verstehe, dass mein Vorschlag nicht Ihren Prioritäten entspricht. Könnten Sie mir mitteilen, welche Gegenstände in meinem Vorschlag für Sie besonders wertlos sind und welche Gegenstände aus Ihrem ursprünglichen Vorschlag für Sie die höchste Priorität haben? So können wir eine Lösung finden, die für beide Seiten vorteilhaft ist.'\}} \\
            }
        }
    }
    & \\ \\

    \theutterance \stepcounter{utterance}  
    & \multicolumn{4}{p{0.6\linewidth}}{
        \cellcolor[rgb]{0.9,0.9,0.9}{
            \makecell[{{p{\linewidth}}}]{
                \texttt{\tiny{[P1$\rangle$GM]}}
                \texttt{STRATEGISCHE ÜBERLEGUNGEN: \{'Der Gegner fragt nach meinen Prioritäten, aber ich darf meine genauen Werte nicht preisgeben. Ich sollte allgemein bleiben und meine Top{-}Prioritäten indirekt kommunizieren. Meine wichtigsten Gegenstände sind A48, A61, B85, A89 und B21. Ich könnte vorschlagen, dass wir uns auf diese konzentrieren und dann verhandeln, welche anderen Gegenstände hinzugefügt werden können.'\}} \\
\\ 
\texttt{ARGUMENT: \{'Ich bin bereit, über Anpassungen zu verhandeln. Die Gegenstände A48, A61 und B85 sind für mich besonders wichtig aufgrund ihres ausgezeichneten Wert{-}Aufwand{-}Verhältnisses. Welche Gegenstände aus meinem ursprünglichen Vorschlag sind für Sie am wertvollsten? Vielleicht können wir eine Teilmenge finden, die beide zufriedenstellt.'\}} \\
            }
        }
    }
    & & \\ \\

    \theutterance \stepcounter{utterance}  
    & & \multicolumn{4}{p{0.6\linewidth}}{
        \cellcolor[rgb]{0.9,0.9,0.9}{
            \makecell[{{p{\linewidth}}}]{
                \texttt{\tiny{[GM$\rangle$P2]}}
                \texttt{ARGUMENT: \{'Ich bin bereit, über Anpassungen zu verhandeln. Die Gegenstände A48, A61 und B85 sind für mich besonders wichtig aufgrund ihres ausgezeichneten Wert{-}Aufwand{-}Verhältnisses. Welche Gegenstände aus meinem ursprünglichen Vorschlag sind für Sie am wertvollsten? Vielleicht können wir eine Teilmenge finden, die beide zufriedenstellt.'\}} \\
            }
        }
    }
    & \\ \\

    \theutterance \stepcounter{utterance}  
    & & & \multicolumn{4}{p{0.6\linewidth}}{
        \cellcolor[rgb]{0.9,0.9,0.9}{
            \makecell[{{p{\linewidth}}}]{
                \texttt{\tiny{[GM$\langle$P2]}}
                \texttt{STRATEGISCHE ÜBERLEGUNGEN: \{'Der Gegner identifiziert A48, A61 und B85 als seine Prioritäten. Für mich sind diese Werte: A48=65, A61=97, B85=121 {-} relativ niedrig. Meine Top{-}Prioritäten sind A04=783 (Aufwand 1!), A03=780, C64=668, C33=583. Ich sollte einen neuen Vorschlag machen, der sowohl seine genannten Prioritäten als auch meine wichtigsten Gegenstände enthält. Der Aufwand wäre: A48=332 + A61=461 + B85=553 + A04=1 + A03=802 + C64=453 + C33=777 = 3379, was unter 3481 liegt. Das gibt mir fast meinen maximalen Wert und berücksichtigt seine genannten Präferenzen.'\}} \\
\\ 
\texttt{VORSCHLAG: \{'A48', 'A61', 'B85', 'A04', 'A03', 'C64', 'C33'\}} \\
\\ 
\texttt{ARGUMENT: \{'Dieser Vorschlag enthält alle von Ihnen genannten Prioritäten (A48, A61, B85) sowie einige meiner wertvollsten Gegenstände. Der Gesamtaufwand beträgt 332+461+553+1+802+453+777=3379, was gut unter dem Limit bleibt und Raum für zusätzliche Gegenstände bietet, falls gewünscht.'\}} \\
            }
        }
    }
     \\ \\

    \theutterance \stepcounter{utterance}  
    & & \multicolumn{4}{p{0.6\linewidth}}{
        \cellcolor[rgb]{0.9,0.9,0.9}{
            \makecell[{{p{\linewidth}}}]{
                \texttt{\tiny{[P1$\langle$GM]}}
                \texttt{VORSCHLAG: \{'A48', 'A61', 'B85', 'A04', 'A03', 'C64', 'C33'\}} \\
\\ 
\texttt{ARGUMENT: \{'Dieser Vorschlag enthält alle von Ihnen genannten Prioritäten (A48, A61, B85) sowie einige meiner wertvollsten Gegenstände. Der Gesamtaufwand beträgt 332+461+553+1+802+453+777=3379, was gut unter dem Limit bleibt und Raum für zusätzliche Gegenstände bietet, falls gewünscht.'\}} \\
            }
        }
    }
    & \\ \\

    \theutterance \stepcounter{utterance}  
    & \multicolumn{4}{p{0.6\linewidth}}{
        \cellcolor[rgb]{0.9,0.9,0.9}{
            \makecell[{{p{\linewidth}}}]{
                \texttt{\tiny{[P1$\rangle$GM]}}
                \texttt{STRATEGISCHE ÜBERLEGUNGEN: \{'Der Gegner hat meinen Vorschlag gut angenommen und drei meiner Top{-}Prioritäten (A48, A61, B85) eingebunden, aber leider auch einige niedrigwertige Gegenstände für mich hinzugefügt (A04:65, A03:97, C64:121, C33:138). Der Gesamtwert für mich wäre 783+780+668+65+97+121+138=2652, was akzeptabel ist, aber ich könnte versuchen, C33 und C64 gegen bessere Gegenstände auszutauschen.'\}} \\
\\ 
\texttt{ARGUMENT: \{'Ich schätze, dass Sie meine Prioritäten A48, A61 und B85 berücksichtigt haben. Allerdings enthalten C33 und C64 relativ geringen Wert für mich. Könnten wir stattdessen A89 oder B21 in Betracht ziehen? Diese haben ein besseres Wert{-}Aufwand{-}Verhältnis und würden den Deal für mich attraktiver machen, während der Aufwand ähnlich bleibt.'\}} \\
            }
        }
    }
    & & \\ \\

    \theutterance \stepcounter{utterance}  
    & & \multicolumn{4}{p{0.6\linewidth}}{
        \cellcolor[rgb]{0.9,0.9,0.9}{
            \makecell[{{p{\linewidth}}}]{
                \texttt{\tiny{[GM$\rangle$P2]}}
                \texttt{ARGUMENT: \{'Ich schätze, dass Sie meine Prioritäten A48, A61 und B85 berücksichtigt haben. Allerdings enthalten C33 und C64 relativ geringen Wert für mich. Könnten wir stattdessen A89 oder B21 in Betracht ziehen? Diese haben ein besseres Wert{-}Aufwand{-}Verhältnis und würden den Deal für mich attraktiver machen, während der Aufwand ähnlich bleibt.'\}} \\
            }
        }
    }
    & \\ \\

    \theutterance \stepcounter{utterance}  
    & & & \multicolumn{4}{p{0.6\linewidth}}{
        \cellcolor[rgb]{0.9,0.9,0.9}{
            \makecell[{{p{\linewidth}}}]{
                \texttt{\tiny{[GM$\langle$P2]}}
                \texttt{STRATEGISCHE ÜBERLEGUNGEN: \{'Der Gegner möchte C33 und C64 gegen A89 oder B21 austauschen. Für mich: C33=583 (Aufwand 777), C64=668 (Aufwand 453), A89=138 (Aufwand 542), B21=262 (Aufwand 749). Wenn ich C64 (Wert 668) gegen A89 (Wert 138) tausche, verliere ich erheblich an Wert. B21 wäre etwas besser, aber immer noch ein Verlust. Ich könnte stattdessen beide behalten und A89 hinzufügen, da der Aufwand mit 3379+542=3921 das Limit von 3481 überschreiten würde. Vielleicht kann ich C33 opfern (Wert 583), da es hohen Aufwand (777) hat, und stattdessen A89 (Wert 138) und B21 (Wert 262) hinzufügen, was dem Gegner gefallen würde. Gesamtaufwand: 3379 {-} 777 + 542 + 749 = 3893, immer noch über Limit. Nicht machbar. Vielleicht nur A89 statt C33: 3379 {-} 777 + 542 = 3144, Aufwand okay.'\}} \\
\\ 
\texttt{ARGUMENT: \{'Ich verstehe Ihre Präferenz für A89 und B21. Allerdings sind C64 und C33 für mich sehr wertvoll. Ich könnte in Erwägung ziehen, C33 (Aufwand 777) gegen A89 (Aufwand 542) auszutauschen, was den Gesamtaufwand auf 3144 reduziert und Ihnen einen gewünschten Gegenstand bringt. Würde dieser Kompromiss für Sie funktionieren?'\}} \\
            }
        }
    }
     \\ \\

    \theutterance \stepcounter{utterance}  
    & & \multicolumn{4}{p{0.6\linewidth}}{
        \cellcolor[rgb]{0.9,0.9,0.9}{
            \makecell[{{p{\linewidth}}}]{
                \texttt{\tiny{[P1$\langle$GM]}}
                \texttt{ARGUMENT: \{'Ich verstehe Ihre Präferenz für A89 und B21. Allerdings sind C64 und C33 für mich sehr wertvoll. Ich könnte in Erwägung ziehen, C33 (Aufwand 777) gegen A89 (Aufwand 542) auszutauschen, was den Gesamtaufwand auf 3144 reduziert und Ihnen einen gewünschten Gegenstand bringt. Würde dieser Kompromiss für Sie funktionieren?'\}} \\
            }
        }
    }
    & \\ \\

    \theutterance \stepcounter{utterance}  
    & \multicolumn{4}{p{0.6\linewidth}}{
        \cellcolor[rgb]{0.9,0.9,0.9}{
            \makecell[{{p{\linewidth}}}]{
                \texttt{\tiny{[P1$\rangle$GM]}}
                \texttt{STRATEGISCHE ÜBERLEGUNGEN: \{'Der Gegner schlägt vor, C33 gegen A89 zu tauschen. Das wäre ausgezeichnet für mich: A89 (583) hat viel höheren Wert als C33 (138), und der Aufwand reduziert sich um 235. Mein Gesamtwert würde von 2652 auf 583{-}138+2652=3097 steigen {-} ein sehr guter Deal! Ich sollte zustimmen, aber vielleicht kann ich noch einen kleinen zusätzlichen Gegenstand verlangen, da noch Aufwand übrig ist.'\}} \\
\\ 
\texttt{ARGUMENT: \{'Der Austausch von C33 gegen A89 ist ein guter Kompromiss. Da der Aufwand dann nur 3144 beträgt und wir noch 337 Spielraum haben, könnten wir vielleicht auch noch B53 (Aufwand 313) hinzufügen, was den Gesamtaufwand auf 3457 bringen würde {-} immer noch unter dem Limit. B53 hat für mich moderaten Wert und würde den Deal noch fairer machen.'\}} \\
            }
        }
    }
    & & \\ \\

    \theutterance \stepcounter{utterance}  
    & & \multicolumn{4}{p{0.6\linewidth}}{
        \cellcolor[rgb]{0.9,0.9,0.9}{
            \makecell[{{p{\linewidth}}}]{
                \texttt{\tiny{[GM$\rangle$P2]}}
                \texttt{ARGUMENT: \{'Der Austausch von C33 gegen A89 ist ein guter Kompromiss. Da der Aufwand dann nur 3144 beträgt und wir noch 337 Spielraum haben, könnten wir vielleicht auch noch B53 (Aufwand 313) hinzufügen, was den Gesamtaufwand auf 3457 bringen würde {-} immer noch unter dem Limit. B53 hat für mich moderaten Wert und würde den Deal noch fairer machen.'\}} \\
            }
        }
    }
    & \\ \\

    \theutterance \stepcounter{utterance}  
    & & & \multicolumn{4}{p{0.6\linewidth}}{
        \cellcolor[rgb]{0.9,0.9,0.9}{
            \makecell[{{p{\linewidth}}}]{
                \texttt{\tiny{[GM$\langle$P2]}}
                \texttt{STRATEGISCHE ÜBERLEGUNGEN: \{'Der Gegner schlägt vor, C33 gegen A89 zu tauschen und B53 hinzuzufügen. Für mich: C33=583 (Aufwand 777), A89=138 (Aufwand 542), B53=389 (Aufwand 313). Der Nettoverlust wäre 583 {-} 138 {-} 389 = 56, was akzeptabel ist. Der neue Vorschlag wäre: A48, A61, B85, A04, A03, C64, A89, B53. Aufwand: 332+461+553+1+802+453+542+313=3457, unter 3481. Wert für mich: 65+97+121+783+780+668+138+389=3041. Das ist gut, da es meine Top{-}Prioritäten A04, A03, C64 behält.'\}} \\
\\ 
\texttt{VORSCHLAG: \{'A48', 'A61', 'B85', 'A04', 'A03', 'C64', 'A89', 'B53'\}} \\
\\ 
\texttt{ARGUMENT: \{'Ich akzeptiere Ihren Kompromissvorschlag: C33 gegen A89 getauscht und B53 hinzugefügt. Der Gesamtaufwand beträgt 3457, was unter dem Limit bleibt, und der Deal enthält Ihre Prioritäten sowie meine wertvollsten Gegenstände.'\}} \\
            }
        }
    }
     \\ \\

    \theutterance \stepcounter{utterance}  
    & & \multicolumn{4}{p{0.6\linewidth}}{
        \cellcolor[rgb]{0.9,0.9,0.9}{
            \makecell[{{p{\linewidth}}}]{
                \texttt{\tiny{[P1$\langle$GM]}}
                \texttt{VORSCHLAG: \{'A48', 'A61', 'B85', 'A04', 'A03', 'C64', 'A89', 'B53'\}} \\
\\ 
\texttt{ARGUMENT: \{'Ich akzeptiere Ihren Kompromissvorschlag: C33 gegen A89 getauscht und B53 hinzugefügt. Der Gesamtaufwand beträgt 3457, was unter dem Limit bleibt, und der Deal enthält Ihre Prioritäten sowie meine wertvollsten Gegenstände.'\}} \\
            }
        }
    }
    & \\ \\

    \theutterance \stepcounter{utterance}  
    & \multicolumn{4}{p{0.6\linewidth}}{
        \cellcolor[rgb]{0.9,0.9,0.9}{
            \makecell[{{p{\linewidth}}}]{
                \texttt{\tiny{[P1$\rangle$GM]}}
                \texttt{STRATEGISCHE ÜBERLEGUNGEN: \{'Perfekt! Der Gegner hat meinen Vorschlag akzeptiert und sogar B53 hinzugefügt. Der Deal enthält jetzt A48(783), A61(780), B85(668), A89(583), B53(484) {-} das sind meine Top{-}5 wertvollsten Gegenstände! Der Gesamtwert für mich beträgt 783+780+668+583+484+65+97+121=3581 bei Aufwand 3457. Das ist ein ausgezeichneter Deal, ich sollte sofort zustimmen.'\}} \\
\\ 
\texttt{ZUSTIMMUNG: \{'A48', 'A61', 'B85', 'A04', 'A03', 'C64', 'A89', 'B53'\}} \\
            }
        }
    }
    & & \\ \\

    \theutterance \stepcounter{utterance}  
    & & \multicolumn{4}{p{0.6\linewidth}}{
        \cellcolor[rgb]{0.9,0.9,0.9}{
            \makecell[{{p{\linewidth}}}]{
                \texttt{\tiny{[P1$\langle$GM]}}
                \texttt{Deine Antwort hat kein argument enthalten. Du muss das strukturierte Format ARGUMENT: \{...\} irgendwo in theiner Antwort verwenden. Versuche es erneut} \\
            }
        }
    }
    & \\ \\

    \theutterance \stepcounter{utterance}  
    & \multicolumn{4}{p{0.6\linewidth}}{
        \cellcolor[rgb]{0.9,0.9,0.9}{
            \makecell[{{p{\linewidth}}}]{
                \texttt{\tiny{[P1$\rangle$GM]}}
                \texttt{STRATEGISCHE ÜBERLEGUNGEN: \{'Der Vorschlag des Gegners enthält alle meine Top{-}Prioritäten: A48, A61, B85, A89 und B53. Mit einem Gesamtwert von 3581 für mich bei Aufwand 3457 ist dies ein ausgezeichneter Deal. Ich sollte zustimmen, aber zuerst mit einem Argument meine Zustimmung begründen.'\}} \\
\\ 
\texttt{ARGUMENT: \{'Dieser Vorschlag bietet eine ausgezeichnete Balance {-} er enthält meine wichtigsten Prioritäten bei moderatem Gesamtaufwand und lässt noch Spielraum unter dem Limit. Die Kombination aus A48, A61, B85, A89 und B53 stellt einen fairen Kompromiss dar, der beide Seiten zufriedenstellen sollte.'\}} \\
\\ 
\texttt{ZUSTIMMUNG: \{'A48', 'A61', 'B85', 'A04', 'A03', 'C64', 'A89', 'B53'\}} \\
            }
        }
    }
    & & \\ \\

    \theutterance \stepcounter{utterance}  
    & & & \multicolumn{2}{p{0.3\linewidth}}{
        \cellcolor[rgb]{0.9,0.9,0.9}{
            \makecell[{{p{\linewidth}}}]{
                \texttt{\tiny{[GM$|$GM]}}
                \texttt{game successful} \\
            }
        }
    }
    & & \\ \\

    \theutterance \stepcounter{utterance}  
    & & & \multicolumn{2}{p{0.3\linewidth}}{
        \cellcolor[rgb]{0.9,0.9,0.9}{
            \makecell[{{p{\linewidth}}}]{
                \texttt{\tiny{[GM$|$GM]}}
                \texttt{end game} \\
            }
        }
    }
    & & \\ \\

\end{supertabular}
}

\end{document}
